%% 第一章 绪论

\chapter{绪论}
\chaptermark{绪论}

\section{研究背景与意义}

无人机在复杂大气环境中的长航时飞行能力是航空工程领域的核心技术挑战之一。
传统固定翼与旋翼飞行器受限于机载能源密度，续航时间通常仅为数十分钟至数小时，
难以满足广域侦察、环境监测和通信中继等长航时任务需求。
自然界中，大型鸟类如安第斯神鹫能够在无扑翼条件下滑翔超过5小时、
覆盖约172公里的距离\cite{xuMigrationSelfpropellingAgent2022}，
其关键在于对大气边界层中热上升气流的巧妙利用。
这种被称为“热滑翔”的行为，本质上是一种在非平衡能量环境中的自适应输运策略：
生物体通过感知局部流场信息，在线决策飞行姿态与航向，
从而将环境中的无序热能转化为定向动能，实现近乎零能耗的远距离迁移。
深入理解这一过程的物理机制与控制逻辑，
对突破无人机续航瓶颈具有重要的理论价值与工程意义。

近年来，活性物质物理学的非平衡统计理论为理解此类现象提供了新的建模框架。
活性物质体系由能够将储存或环境能量转化为定向运动的个体组成，
其集体动力学行为展现出超越平衡态统计物理的丰富现象
\cite{tonerHydrodynamicsPhasesFlocks2005}。
将无人机滑翔纳入活性物质理论的分析范式，
能够从统计物理视角揭示滑翔系统在湍流能量场中的平均行为与随机响应规律
\cite{ramaswamyMechanicsStatisticsActive2010}，
为理解其滑翔策略、能耗优化和环境交互提供新的建模思路。
活性物质理论的核心优势在于：它将个体运动、环境能量场和群体统计行为统一在
非平衡态热力学的框架下，从而能够自然地处理大气湍流所固有的
多尺度、非定常和随机性特征。通过建立无人机滑翔与活性颗粒运动之间的映射关系，
可以实现从微观活性系统到宏观空气动力学系统的跨尺度理论统一，
这不仅拓展了活性物质理论的应用边界，也为非平衡系统的统计建模提供了物理依据。

与此同时，深度强化学习的引入为无人机滑翔的非平衡能量优化提供了全新的技术路径。
传统的滑翔轨迹规划方法依赖对风场结构的先验建模与显式求解，
在真实大气的高维非定常特征面前面临维度灾难与模型失配的双重困境。
深度强化学习以其无模型、数据驱动的特性，能够通过与环境的持续交互
自主学习最优控制策略，从而绕过复杂的物理建模过程。
张进等的研究表明，深度强化学习在非线性流体系统中能够自动识别能量传递规律，
实现高效节能控制\cite{ZhangJinQiangHuaXueXiFangFaZaiYiXingPaiDongShiYanZhongDeYingYong2023}；
Colabrese等利用强化学习优化微型泳者在湍流中的路径选择，
其结果显示智能体可学习出最优能量轨迹，
这与滑翔无人机在能耗最小化任务中的策略学习具有高度相似性
\cite{colabreseSmartInertialParticles2018a}。
基于机器学习的能耗优化模型能在湍流环境下实现自学习、
自适应与能量分配的动态调整：
结合活性物质的能量输运模型，
强化学习算法能够模拟滑翔无人机的非平衡能量利用行为，
实现对复杂大气能场的智能感知与高效响应，
为发展新一代自主节能飞行器奠定理论与算法基础。

综上所述，无人机滑翔不仅是一个空气动力学问题，
更是一个典型的非平衡能量系统优化问题。
它结合了活性物质物理、湍流动力学与人工智能控制三大领域，
为研究多尺度环境中的能量利用与迁移策略提供了一个理想模型体系。
滑翔无人机在长航时任务、节能飞行及自主控制领域具有广泛应用前景：
通过基于活性物质理论的能量优化模型，
可显著提升无人机在风能环境中的续航性能与任务效率。
这一研究思路可进一步推广至仿生飞行器、能源俘获系统
和智能群体飞行控制等领域，
对推动绿色航空与智能空域管理具有重要意义。
\section{国内外研究现状}

\subsection{无人机滑翔研究}

无人机滑翔研究源于对鸟类滑翔机理的长期探索，其核心目标是在无动力或弱动力条件下
最大化飞行航程与续航时间。这一问题的本质在于：如何在具有强时空变异性的
大气能量场中，以最小的自身能耗实现最远的迁移距离。
早期工作主要集中于热气流滑翔的物理基础与观测特征。
Woodward从气象与滑翔机飞行角度建立了孤立热气流的结构模型，指出热泡内部存在
明显的上升核与径向流动，周围伴随补偿性下沉气流，这一环状涡结构
为后续热滑翔建模奠定了流体力学基础\cite{woodwardTheoryThermalSoaring1956}。
Cone则通过系统的野外观测，总结了鸟类“盘旋—滑翔”交替飞行的典型模式，
强调翼载荷、盘旋半径与热气流强度之间的定量关系——翼载荷越小、
盘旋半径与热气流核心半径越匹配，能量获取效率越高，
这些结论至今仍是无人机热滑翔设计的重要依据\cite{coneThermalSoaringBirds1962}。
在飞行性能优化方面，MacCready提出的最优空速选择理论为跨热区滑翔速度规划
提供了经典框架：该理论将热气流平均上升速率与滑翔机下沉极曲线结合，
推导出使跨区平均速度最大化的最优飞行速度，
并被证明同样适用于鸟类与滑翔机\cite{maccreadyOptimumAirspeedSelector1958}。

随着高精度定位与轨迹数据的积累，研究逐步从定性描述转向定量分析。
Ákos等利用GPS轨迹对比了隼类与人类滑翔伞飞行，发现鸟类在跨区滑翔阶段
呈现出近似“MacCready规律”的速度选择行为，
且隼类的速度调整比人类滑翔伞更为精细和及时，
说明自然进化与工程优化在策略层面具有一致性，
同时暗示鸟类可能拥有比现有工程模型更优的感知-决策闭环机制
\cite{akosComparingBirdHuman2008}。
Shannon等进一步结合气象模型，揭示了鹈鹕飞行高度与热气流深度、
强度之间的对应关系，发现鹈鹕倾向于在热气流有效深度
（即上升速率超过下沉速率的区域）的上半部分盘旋，
表明鸟类能够根据边界层结构实时调整滑翔包线，
这为无人机感知与决策算法提供了重要的生物学启示
\cite{shannonAmericanWhitePelican2002a}。

在工程领域，研究重点逐渐转向动态滑翔与最优轨迹规划。
Zhao建立了滑翔机在梯度风场中的动态滑翔最优控制模型，
推导了在已知风场条件下能量最优的周期轨道的必要条件，
揭示了动态滑翔的能量增益来源于飞行器在风梯度层间的
反复穿越所实现的净风能提取\cite{zhaoOptimalPatternsGlider2004}。
Sachs等将最优控制理论系统应用于动态滑翔轨迹优化，
证明了在理想剪切风场中，动态滑翔可实现的最小风梯度需求
与无人机气动效率之间存在明确的函数关系\cite{sachsApplicationOptimalControl2005a}。
Liu等进一步研究了动态滑翔的最优机动模式，发现除经典的
“闭环O型”轨迹外，还存在“开口S型”等多样化机动形式，
不同模式适用于不同的风速条件和任务需求\cite{liuOptimalPatternsDynamic2017}。
Cui等针对滑翔机在热气流中的能量获取问题，
建立了包含热气流模型和飞行器动力学的耦合仿真框架，
定量评估了不同盘旋策略的能量收益差异\cite{cuiStudyGliderSoaring2023}。

近年来，随着人工智能技术的快速发展，深度强化学习为无人机滑翔提供了
全新的技术路径。与传统方法相比，强化学习无需对风场结构进行显式建模，
而是通过与环境的持续交互自主学习控制策略。
Reddy等率先将强化学习应用于滑翔问题，使用Q-learning算法训练
滑翔机在湍流热气流场中自主决策，发现智能体能够学习到类似于
鸟类“盘旋-滑翔”的行为模式，且其能耗显著低于预设的基准策略，
首次验证了无模型方法在热滑翔任务中的可行性
\cite{reddyGliderSoaringReinforcement2018}。
在此基础上，Reddy等进一步证明了在未知湍流环境中，
强化学习智能体能够学会利用局部感知信息进行在线路径规划，
在无任何先验风场知识的条件下实现了鲁棒的自主滑翔，
这一发现打破了长期以来认为滑翔控制必须依赖精确风场模型的工程范式
\cite{reddyLearningSoarTurbulent2016}。
Cui等将深度强化学习拓展至三维风场环境，
通过引入高度通道的独立感知与决策机制，
实现了水平路径与垂直剖面的联合优化，
使滑翔机能够自主选择最佳飞行高度层以匹配热气流分布
\cite{cuiStudyGliderSoaring2023}。
Lawrance在其博士论文中系统研究了自主滑翔飞行的理论与工程实现，
构建了从风场建模、轨迹规划到在线决策的完整框架，
并讨论了将强化学习嵌入滑翔控制回路的可行方案
\cite{lawranceAutonomousSoaringFlight2011}；
其风能路径规划工作进一步提出了基于风场能量地图的
长航程滑翔策略\cite{lawranceWindEnergyBased2009}。

Edwards团队在工程实现层面进行了开创性探索。他们开发了自主滑翔
飞控系统，在北卡罗来纳州进行了实际飞行验证，
证明自主热滑翔可以在真实大气条件下实现净能量增益
\cite{edwardsAutonomousSoaringMontague2010}；
其后续工作详细阐述了机载传感器的数据融合、
热气流中心定位算法和飞行控制律的具体实现细节，
为从仿真到实飞的工程化过渡提供了完整的技术路线
\cite{edwardsImplementationDetailsFlight2008}。
He等进一步利用深度强化学习在非平稳风场中训练滑翔无人机，
发现智能体能够自发学习出能量最优的“跨区直飞—热区盘旋”策略，
在完全未知的风场分布中实现了正的能量净收益，
证明了DRL方法在面对真实大气不确定性和非平稳性时的适应能力
\cite{heEnergypositiveSoaringUsing2024}。

在动态滑翔的深度强化学习前沿，Flato等受秃鹫在极端强风与湍流环境中的
飞行行为启发，探索了自主热气流滑翔的鲁棒控制原理。
研究表明，经过深度强化学习训练后，神经网络的隐层表征会自发形成
专门应对不同气流扰动模式的“功能聚类”——某些神经元群体
对剪切湍流敏感，另一些则偏好热泡上升气流，
这种隐式编码的多模态环境感知能力极大提升了无人机
在强随机气象条件下的生存与能量捕获鲁棒性
\cite{flatoThermalSoaringDRL2024}。
Chen等证明了在非定常剪切流中，无人机动态滑翔无需依赖
预先规划的“周期级”参考轨迹。通过深度强化学习，
智能体仅凭局部观测就能自然涌现出“步进级”控制律——
每一步决策只依赖当前状态和短期记忆而无需全局轨迹信息，
这种去中心化的控制架构在“动态能量收集”与“定向目标滑翔”之间
实现了自适应的实时权衡，为解释信天翁等海鸟在极端海况下
高效动态滑翔的神经控制机制提供了计算层面的有力证据
\cite{chenDynamicSoaringDRL2026}。

在多算法对比与工程验证方面，Adamski等将强化学习推向了高保真度应用。
该研究结合高保真CFD计算与涡格法（VLM），
利用深度确定性策略梯度（DDPG）算法在六自由度无人机模型上
成功训练了闭环路径跟随与无制导风场能量搜索智能体。
结果表明，即便在考虑全机气动力非线性、舵面饱和
和传感器噪声等真实工程约束的情况下，
DRL策略依然能够有效地在未知风场中定位并利用上升气流，
验证了强化学习在强非线性气动环境中的工程可行性
\cite{adamskiTowardsDevelopmentDynamic2023}。
Park等系统性地评估和对比了TD3、SAC、PPO等多种
先进深度强化学习算法在动态滑翔任务中的表现。
研究发现，基于最大熵框架的SAC算法在探索效率
和策略稳定性方面优于确定性策略梯度方法，
而PPO在样本效率上具有优势。各算法训练的智能体均能有效学习滑翔机动，
在完全未知的流场中最大化利用风梯度提取能量并最小化能量损失
\cite{parkApplicationReinforcementLearning2025}。

在控制架构与感知方法创新方面，Kim与Perez提出了针对无人机动态滑翔的
鲁棒神经控制架构。该架构采用在线自适应神经网络，
能够在飞行过程中根据实时感知的风场变化持续更新控制策略，
验证了在具有强随机性与非定常特性的风梯度场中，
利用神经网络进行自主能量获取和轨迹控制的可靠性
\cite{kimRobustNeurocontrolAutonomous2022}。
Zhao等将Dryden连续湍流模型与Gedeon热气流模型耦合，
利用TD3算法探索了滑翔机在完全非稳态随机扰动气流中的滑翔策略，
证明了即使在叠加了连续随机湍流脉动的恶劣气象条件下，
智能体依然能够通过自适应调整盘旋中心和盘旋半径来实现能量正收益，
这为无人机在真实大气环境中的全天候自主滑翔提供了算法基础
\cite{zhaoEnergyHarvestingStrategy2023}。

在架构创新与方法拓展方面，Abozeid等对强化学习在动态滑翔中的应用潜力
进行了全面的理论评估与机制剖析。研究者对比了策略梯度、
值函数和Actor-Critic三类算法的收敛特性与泛化性能，
指出强化学习方法有望取代传统的极值搜索算法，
在无需显式风场建模的条件下实现复杂风场环境中的高效轨迹规划，
为该领域的工程化应用提供了系统性的理论依据
\cite{abozeidComprehensiveAssessmentPotential2023}。
Gall等针对真实环境中的局部感知瓶颈问题，
引入了时间卷积网络（TCN）进行端到端的热上升气流检测与风场结构估计。
与传统基于卡尔曼滤波或批处理优化的方法不同，
TCN能够通过扩张卷积在长时序列中捕获多时间尺度的流场特征，
从而在传感器数据稀疏和噪声干扰严重的条件下
显著提升滑翔机对复杂环境的预判与适应能力
\cite{gallEndtoEndThermal2024}。
Schimpf等将多智能体强化学习（MARL）应用于无人机集群热滑翔，
通过构建无人机群体间的协作通信机制与分布式策略学习框架，
验证了智能体之间通过局部信息交互，
能够自发形成空间分散的协同搜索队形，
显著缩小搜索盲区并提升集群整体的能量捕获效率
\cite{schimpfMultiAgentReinforcementLearning2021}。
Liu与Li进一步提出了一种融合热成像视觉信息
与强化学习的多模态滑翔框架，
利用机载红外热像仪提供的空间温度分布作为额外观测通道，
使智能体能够提前感知并精确识别热气流边界和核心位置，
从而将滑翔轨迹的规划从“被动响应”提升为“主动预判”，
显著改善了在未知热气流场中的轨迹生成质量
\cite{liuThermalImageGuided2026}。

上述研究表明，深度强化学习无需依赖先验风场模型或显式轨迹规划，
仅凭局部环境感知即可使无人机在非定常剪切流中自主涌现出
高效的能量捕获与目标导向滑翔策略。从早期基于Q-learning的简单表格策略，
到Actor-Critic架构下的连续动作空间优化，
再到多模态感知融合和多智能体协同，
强化学习正逐步从仿真验证走向工程实用，
为理解生物滑翔机制和工程化无人机节能飞行提供了理论支撑与算法保障。
\subsection{活性物质非平衡输运研究}

活性物质输运的理论根源可以追溯到Zermelo于1931年提出的经典导航问题：
在已知风场分布的条件下，飞行器如何选择航向以最小化到达时间
\cite{zermeloUberNavigationsproblemBei1931}。
该问题的数学本质是在非均匀矢量场中求解时变最优控制问题，
其最优解满足带漂移项的Hamilton-Jacobi-Bellman方程。
Mureşan对该问题进行了系统性的数学分析，揭示了在定常流场中
最优路径与Hamilton-Jacobi方程的内在联系，
并给出了在不同流场拓扑结构下最优解的存在性与唯一性条件
\cite{muresanZermelosNavigationProblem2014}。
然而，经典Zermelo框架依赖于对全局流场的精确先验知识，
当流场具有湍流特征时，速度场在空间和时间上均表现出
宽频谱的随机脉动，使得最优解的解析构造不再可能。
Reddy等在湍流热气流环境中训练滑翔智能体，
发现其习得的策略能够在无模型条件下实现路径随机性的自适应利用——
智能体学会在湍流脉动有利时顺流加速、不利时主动规避，
这种在线概率推理能力突破了经典Zermelo框架对确定性风场模型的依赖
\cite{reddyLearningSoarTurbulent2016}。

在微观尺度上，Colabrese等提出了“智能惯性颗粒”的概念，
将活性物质理论与强化学习相结合，通过训练点状颗粒
在二维湍流场中完成点对点迁移任务。
研究发现，经过训练的智能颗粒能够学会利用流场中的涡结构——
主动靠近涡的边缘以借助切向速度加速前行，
其能量消耗显著低于沿直线运动的基准方案，
且节省的能耗比例随湍流强度的增大而提高，
说明湍流在提供扰动的同时间也蕴含着可供利用的有序结构
\cite{colabreseSmartInertialParticles2018a}。
这一发现为后续研究奠定了重要的方法论基础：
将活性颗粒的导航问题形式化为强化学习中的序贯决策问题，
使智能体能够通过与流场环境的反复交互
自动发现并利用其中的能量输运通道。
Xu等将这一框架推广至瑞利-伯纳德热对流系统，
在该系统中大尺度环流与热羽流共同构成了
远比二维各向同性湍流复杂的多尺度流动结构。
研究验证了即便在如此复杂的湍流环境中，
强化学习策略仍然能够有效降低迁移能耗，
并发现将温度场信息作为额外观测通道可以显著提升训练效率，
表明多物理场感知对于复杂环境中的智能导航具有重要价值
\cite{xuMigrationSelfpropellingAgent2022}。
在此基础上，Xu等进一步研究了大宽高比腔体中的长距离迁移问题，
结果表明智能体能够学习并利用大尺度环流结构作为“高速公路”，
在整个腔体宽度上实现长距离的高效定向输送，
且节能效果随腔体宽高比的增大而增强，
反映了流场有序结构的空间延展性对迁移效率的促进作用
\cite{xuLongdistanceMigrationMinimal2023}。

Monthiller等从浮游生物导航的生物学视角出发，
提出了“湍流冲浪”策略——生物体或智能体只需通过被动感知
局部流场梯度信息（如速度的旋度或应变率），
就能判断有利的流动方向并实现顺流而行，
其能耗在统计意义上接近理论最优值。
该策略的美妙之处在于其极低的感知与计算需求：
不需要对全局流场的记忆或预测，仅凭瞬时的局部信息即可做出
接近最优的导航决策，这为资源受限的微型机器人提供了
极具吸引力的控制范式\cite{monthillerSurfingTurbulenceStrategy2022}。
Monthiller的博士论文进一步从力学机理层面系统阐述了
浮游生物在湍流中的迁移策略，揭示了感知与响应之间
存在一个由流体力学时间尺度决定的特征耦合时间——
感知过快会引入噪声，过慢则会错失有利流动窗口，
最优感知频率恰好与流场的能量含能尺度频率相匹配
\cite{monthillerMechanisticApproachPlankton2022a}。
Mousavi等从生态学视角研究了浮游动物在湍流中的高效生存策略，
发现在捕食压力与能耗约束的双重选择压力下，
最优策略呈现出“间歇性主动运动”与“被动随流漂移”交替的
burst-and-coast特征，这种节律性行为在降低代谢消耗的同时
保持了足够的环境搜索效率\cite{mousaviEfficientSurvivalStrategy2024a}。
Calascibetta等从最优跟踪视角出发，分析了湍流中一个活性颗粒
追踪另一个移动目标的最优策略，
发现当追踪者学会利用流场中的相干结构时，
其追踪能耗可降低至直线追踪的40\%以下，
且追踪成功率反而因轨迹的流线顺应性而得到提升
\cite{calascibettaOptimalTrackingStrategies2023}。

上述研究表明，无论是宏观尺度的无人机滑翔，还是微观尺度的颗粒输运，
活性物质体系的迁移行为均展现出“顺应而非对抗”流场的共性规律——
即利用流场中的有序结构（涡、剪切层、环流等）作为能量载体，
以最小的自身能耗实现最大范围的空间迁移。
然而，当前研究多在微尺度系统中展开，流动以低雷诺数层流或弱湍流为主，
面向宏观飞行器场景的高雷诺数湍流条件下的能耗优化建模
与分析仍相对不足，亟需建立跨尺度、跨雷诺数的统一理论框架。

近年来，强化学习在活性物质领域的应用持续深化，
研究范畴从个体导航策略拓展至群体动力学调控和系统边界优化。
Cai等的综述系统梳理了强化学习在活性物质物理中的多层次应用，
从个体自推进微粒在复杂流场中的最优路径寻优，
到通过外部智能体调控活性浴的集体行为以实现特定功能，
再到利用强化学习优化活性系统的非平衡态热力学边界，
勾勒出该交叉领域从微观到宏观、从单智能体到多智能体的完整研究图谱
\cite{caiRLActiveMatter2025}。
Nasiri等首次使用基于策略梯度的深度强化学习方法，
在无需人为设计分段奖励函数（reward shaping）的条件下，
为微观活性粒子在包含随机障碍物和局部流场畸变的
复杂崎岖环境中规划出渐进最优的迁移路径。
这一工作的重要性在于证明了端到端强化学习可以取代
传统需要大量领域专家知识的分段导引策略，
使活性颗粒导航系统具备更强的泛化能力和环境适应性
\cite{nasiriRLActiveNavigation2022}。
同一团队进一步揭示了强化学习驱动的智能微粒
在长期训练中会自动演化出类似于趋化性细菌的
“跑-翻滚”（run-and-tumble）运动节律——
沿梯度方向直线游动（跑）与随机重定向（翻滚）交替进行，
且两者的持续时间比随环境梯度强度的变化而自适应调节，
这种从零开始自发涌现的生物拟态行为深刻揭示了
物理约束与优化目标在策略空间中的收敛方向
\cite{nasiriSmartParticlesBacterial2024}。
在热对流系统中，Zhou等将强化学习应用于湍流热对流的边界控制问题，
展示了强化学习智能体通过非线性动态调节壁面温度的空间分布和时序，
能够主动重构对流涡的结构和输运效率，
最终打破了经典Kraichnan尺度律所预测的热传导极限，
实现了高达38.5\%的物理热传导效率提升，
这一结果彰显了强化学习在非平衡系统主动调控中的巨大潜力
\cite{zhouDeepReinforcementLearning2025}。

在导航策略优化方面，Putzke与Stark研究了微型游动体在信息严重缺失条件下的
智能导航极限问题。在仅能感知相对目标的方向和距离、
且叠加显著热噪声（布朗运动）的欠观测条件下，
研究证明即便在具有强剪切的泊肃叶流及非线性旋涡等复杂流场中，
Q-learning算法依然能够使微观颗粒学得收敛于耗时最短的迁移路径。
该结果从信息论角度确立了一个重要下界：
智能导航所需的最小感知信息量可以低至仅两个标量（方向与距离），
这对于设计传感能力极为有限的微型医疗机器人具有重要指导意义
\cite{putzkeOptimalNavigationSmart2023}。
Zhang等在高保真计算流体力学（CFD）环境中训练仿生鱼类游动体，
在综合考虑流体动力、身体变形和涡流相互作用的完整物理模拟中，
通过在强化学习奖励函数中引入与瞬时功率成正比的能耗惩罚项，
观察到智能体自发淘汰了耗能巨大但推进效率低下的高频连续扑动，
演化出与真实鱼类高度一致的“爆发-滑行”（burst-and-coast）节能步态——
短时高强度加速后被动滑行以回收能量。
更深入的流场分析表明，滑行阶段的体态保持和尾迹利用
是能耗降低的关键物理机制
\cite{zhangReinforcementLearningRobust2025}。
Lai等面向低雷诺数下的三连杆微型游动体，
系统设计了“速度优先”与“能耗感知”两种差异化的强化学习奖励函数，
清晰地展示了奖励信号的设计如何从根本上决定游动体变形步态的分化。
速度优先模式下游动体倾向于高频大振幅的非对称变形以最大化瞬时速度，
而能耗感知模式下则演化出低频、小振幅且沿流线方向对齐的
对称变形模式以最小化形状阻力，
这一对比实验为活性游动体的“形态-步态-能效”关系研究
提供了可靠的参数化分析框架
\cite{laiNavigationThreeLinkMicroswimmer2025}。
Chakraborty等面向靶向药物递送等医疗应用场景，
研究了微型游动体在泊肃叶流中的逆流寻的导航问题。
智能体能够根据感知到的局部流速和剪切率的变化，
连续且精准地在“摆动”（wiggling）与“翻滚”（tumbling）两种
基本运动模式之间切换——高剪切区域采用摆动以克服流体阻力，
低剪切区域切换为翻滚以调整朝向，
这种模态自适应切换策略实现了从任意起始位置向目标区域的高效逆流迁移
\cite{chakrabortySmartNavigationMicroswimmers2025}。

在流场尺度与建模方法方面，Parfenyev针对具有逆能量级联特性的
充分发展二维湍流重新求解了Zermelo导航问题。
在充分发展湍流中，能量从注入尺度向大尺度转移形成有序的涡街结构，
这些大尺度相干结构占据了流场动能的主要部分。
该研究得出一个重要结论：仅利用滤除高频脉动后的大尺度粗粒化流场特征
进行训练的策略，能够完美泛化至包含全尺度复杂脉动的真实流场中——
因为控制迁移能耗的物理过程主要由大尺度流动结构决定，
高频小尺度脉动对平均迁移效率的影响在统计上相互抵消
\cite{parfenyevOptimalNavigationTwoDimensional2025}。
Tonti与Vinuesa针对无人机在三维城市建筑群湍流中导航的工程难题——
该场景涉及建筑物尾流中的脱落涡、街道峡谷中的回流区
以及屋顶剪切层的强湍流混合等复杂流动现象——
提出了将Proximal Policy Optimization (PPO)
与GTrXL（基于Transformer的时序特征提取器）
或LSTM（长短期记忆网络）时序记忆网络相结合的深度强化学习架构。
时序记忆模块有效地克服了局部流场感知在单体观测中
带来的“视野盲区”和“部分可观测马尔可夫决策过程”困境，
使智能体能够推断当前所处的高维湍流拓扑结构，
最终将导航成功率从传统方法的不足60\%提升至98.7\%，
坠毁率降至0.1\%
\cite{tontiNavigationThreeDimensionalUrban2025}。
Mousavi等从时间尺度分析的角度，
深入对比了短时间视界与长时间视界下微观活性游动体的
强化学习导航差异。研究发现，短视界策略倾向于最大化
单步推进效率而忽略长周期流场结构（如大尺度环流），
可在弱湍流条件下取得较好的局部性能但缺乏全局视野；
长视界策略则主动规划绕行路径以利用大尺度对流的输运优势，
虽然在即时奖励上不如短视策略，但总迁移时间和能耗更优。
该研究为在复杂湍流中平衡即时贪婪效益与长期全局最优
提供了具体的算法配置建议——对于大尺度流场占优的环境，
应适当增大折扣因子以赋予智能体更长远的规划视野
\cite{mousaviShortTermLongTerm2025}。
\subsection{强化学习在流体控制中的应用}

近年来，深度强化学习在流体力学领域取得了显著进展，
研究范围从宏观集群运动控制延伸至微观游动策略优化和主动流动控制。
这一交叉领域的核心科学问题是：
如何设计通用的学习框架，使智能体能够在
高维、非线性和非定常的流体动力环境中
自主发现最优的控制策略。

在仿生集群游动方面，Gazzola等率先利用强化学习训练鱼群
在考虑完整流体动力相互作用的模拟环境中自主游动，
发现智能体能够在不依赖任何预设编队规则的情况下，
自发学习出有利于减小群体平均阻力的菱形和交错型编队构型，
且编队构型的拓扑结构随群体规模和游动速度的变化而自适应调整
\cite{gazzolaLearningSchoolPresence2016}。
Verma等将深度强化学习应用于鱼群协同游动问题的更深入研究，
结果表明跟随智能体不仅能够精准定位前导鱼尾迹中的
低压力区和高动量区以实现“借力而行”，
更令人惊讶的是，整个鱼群系统通过强化学习
涌现出了类似真实鱼群中观察到的“领导-跟随”层级结构，
而这一结构并非由奖励函数显式编码，
而是智能体在流体动力约束下最大化自身推进效率的自然结果
\cite{vermaEfficientCollectiveSwimming2018}。

在微尺度主动输运领域，Muiños-Landin等将强化学习引入人工微泳者系统，
不仅在数值模拟中验证了强化学习策略的有效性，
更进一步在真实的人工微泳者物理实验平台上实现了
强化学习策略的在线部署与验证，
克服了仿真到现实（sim-to-real）迁移中的物理建模误差和传感噪声等挑战，
为强化学习在微尺度物理系统中的实际应用提供了关键的实验证据
\cite{muinos-landinReinforcementLearningArtificial2021}。
Mirzakhanloo等提出了“主动隐身”的前沿概念——
通过强化学习优化微型泳者的游动策略，
使其在斯托克斯流中以特定波形和相位运动，
使其诱导的远场流动在时间平均意义上抵消为零，
从而对周围流体环境产生最小的扰动。
这一工作不仅拓展了活性物质控制的物理边界，
也为隐蔽性微型水下航行器的设计提供了全新的思路
\cite{mirzakhanlooActiveCloakingStokes2020}。
Monderkamp等从统计物理角度出发，证明了活性布朗颗粒
通过表格型Q-learning可以在无需任何先验模型的情况下，
学习到在复杂随机势能景观中高效导航的策略，
且习得的策略展现出对未见过的势能分布和环境参数的
强大泛化能力，这为理解生物细胞在复杂细胞外基质中
定向迁移的物理机制提供了计算模型支撑
\cite{monderkampActiveParticlesUsing2022}。
Piro等从能量学角度系统分析了微泳者身体形状
对其在湍流中导航能耗的定量影响，
揭示了细长体形态和曲率分布等几何参数
与策略最优性之间的非单调耦合关系——
过长或过短的身体都会增加在湍流中保持朝向的代谢代价，
存在一个与流场相关长度尺度相匹配的最优体形比
\cite{piroEnergeticCostMicroswimmer2024}。

以上研究表明，强化学习在处理非定常、强非线性流体系统的
最优控制问题上具有独特的优势。其无模型、数据驱动的特性，
使其特别适用于难以建立精确数学模型的湍流环境——
智能体能够通过反复试错，从高维感官输入中提取与任务相关的
低维流场特征，并在策略空间中自主搜索全局最优解。
然而，现有研究主要集中在微尺度系统和低雷诺数流场中，
流动拓扑相对简单，非线性程度和自由度数量有限。
面向无人机尺度的宏观滑翔能耗优化研究尚处于起步阶段：
宏观飞行器面临的高雷诺数湍流具有更宽的频谱范围、
更强的非线性耦合和更复杂的边界层分离现象，
对强化学习的样本效率、策略泛化性和实时在线决策能力
均提出了更高的要求，亟需开展系统的数值模拟与理论分析。

在宏观尺度的主动流动控制（Active Flow Control, AFC）领域，
深度强化学习同样展现出突破性的控制能力，
正在深刻改变流体力学中经典钝体绕流控制的研究范式。
传统AFC方法——如周期性激励、开环吹吸气、
基于简化模型的反馈控制等——在面对高雷诺数下的
宽频湍流脉动和强非线性流动分离时往往力不从心。
深度强化学习则通过学习环境动态与最优控制动作之间的
端到端映射，为这一经典难题提供了全新的求解思路。

Moslem等的综述系统回顾了深度强化学习在钝体主动流动控制中的应用全景，
从控制目标（减阻、抑制振动、增强混合）到算法选择
（DDPG、PPO、SAC及其变体），全面论证了DRL在
多频率涡流同步抑制、实时自适应参数调节和
鲁棒性方面的显著先进性，
并指出了制约该领域发展的关键瓶颈——高保真CFD环境下
的巨大样本需求和训练时间成本
\cite{moslemDeepReinforcementLearning2025}。
Montalà等将DRL控制器与GPU加速的浸入边界法求解器深度耦合，
在三维NACA0012翼型在大迎角条件下的全湍流分离流中，
实现了闭环自主气动优化。智能体通过调节分布于翼面的
多个独立作动器的幅值和相位，
重新组织了吸力面的分离剪切层结构，
使其再附点显著前移，从而实现了升力的恢复与阻力的降低
\cite{montalADRLFlowSeparation2025}。
Jia等在经典圆柱绕流基准算例中运用DRL操纵一对合成射流作动器，
成功将平均阻力系数降低了49\%，并彻底消除了交替脱落的卡门涡街——
智能体学会以特定的频率和相位差驱动上下表面射流，
在尾流中产生了对称的“阻塞效应”使分离泡对称化，
实现了混沌尾迹向流线型对称尾迹的拓扑转变
\cite{jiaActiveFlowControl2025}。
Ren等将DRL应用于弹性支撑圆柱体的自转控制以抑制涡致振动（VIV），
在仅提供圆柱位移和速度等少量运动学观测的情况下，
智能体学会了根据实时振动相位主动驱动圆柱旋转
以破坏尾流涡的展向相关性和锁定（lock-in）效应，
将振幅降低了99.6\%，同时显著增强了流固耦合系统的模态稳定性，
为深海立管和大跨度桥梁的涡振抑制提供了智能化解决方案
\cite{renDRLVIV2024}。
Gunnarson等从更基础的流体力学视角出发，
系统验证了强化学习智能体在时变非定常双涡流环境中的路径规划能力。
在该基准流场中，点涡的位置和强度随时间按预设规律变化，
智能体需要在没有任何全局流场信息的情况下，
仅凭局部速度感知规划从起点到终点的最优轨迹。
结果表明经过训练的智能体能够捕获涡流外围的有利速度区域
进行“乘波而行”式的高效迁移，
其轨迹与利用全局信息的理论最优解高度吻合
\cite{gunnarsonNavigationVortical2021}。

在更广泛的流动控制场景中，Jia Wang等将强化学习从标准圆柱
拓展至椭圆到平板等横截面长宽比连续变化的变截面钝体，
探索了单一控制策略在多种气动外形间的迁移泛化能力。
通过PPO算法精准控制分布于钝体表面的多个合成射流的质量通量，
使射流动量与尾流剪切层发生非线性耦合，
在截面长宽比从1（圆柱）到无穷大（平板）的宽参数空间内
均有效抑制了卡门涡街的非对称交替脱落，
实现了橫跨不同钝体形态的大幅气动减阻
\cite{jiaWangDeepReinforcementLearning2024}。
Montalà等以题为“Towards Active Flow Control Strategies Through
Deep Reinforcement Learning”的研究，
标志着基于数据驱动的主动流动控制从二维简化流场
向三维全流域真实物理控制的跨越。
在雷诺数为100的三维圆柱绕流中应用深度强化学习框架，
智能体成功学习到沿展向非均匀分布的高级作动策略——
部分展向位置采用同相位吹吸以抑制主模态脱落，
另一部分位置采用反相位激励以破坏展向涡的相干性，
实现了9.32\%的直接减阻与高达78.4\%的升力振荡下降，
这一多模态组合策略体现了三维流场控制的丰富物理内涵
\cite{montalaTowardsActiveFlowControl2024}。
Liu等系统性地评估了翼型主动流动控制中超参数配置
对最终流体力学性能的定量影响，
聚焦于控制更新间隔（control update interval）
和状态观测空间配置两个关键超参数。
研究证实，合理的控制更新间隔——既不太短以致引入高频噪声，
也不太长以致无法及时响应流场变化——
能为智能体引入更广泛的作动频率覆盖范围，
使其能够与从低频大尺度分离泡到高频剪切层
Kelvin-Helmholtz不稳定性等更宽频带的流场脱落结构
进行有效交互并实施定向抑制
\cite{liuReinforcementLearningClosedLoop2025}。
Xu等将强化学习智能体置于一个横跨从弱湍流到硬湍流
（瑞利数$Ra$从$10^6$到$10^{11}$，跨越五个数量级）
的极端瑞利-伯纳德热对流系统中进行点对点导航训练。
在这一宽广的参数范围内，系统的流场拓扑经历了从
稳态层流到周期性振荡、再到完全发展的时空混沌
和终极区湍流的完整转捩路径。
研究揭示了随着雷利数增加和大尺度环流结构的拓扑破碎，
智能体寻找低能耗迁移通道的策略也随之
发生本质性的范式转变——在中等$Ra$下，
智能体的策略表现为“强行突破大尺度环流的输运势垒”，
即利用自推进力克服对流环的阻碍；
而在极高$Ra$的硬湍流状态下，随着大尺度环流瓦解为
大量瞬态爆发的随机热羽流，
智能体自然演变为“利用瞬态上升羽流的间歇性跃迁”策略——
等待并捕捉局部上升热羽流以借力跃升，
随后在羽流间期的最低能耗状态下滑行，
这种从“对抗流场”到“顺应流场”的策略转变
揭示了活性物质在极端非平衡环境中节能导航的普适物理机制
\cite{xuLearningTraverseConvective2026}。
Nair与Goza研究了受鹰鸮等鸟类飞羽变刚度机制启发的
变刚度后缘襟翼，通过强化学习同时优化襟翼的结构刚度分布
和动态偏转策略，实现了对不同来流条件和气动载荷状态的
自适应混合流场控制。在这一流固耦合控制框架中，
结构参数（刚度分布）决定了被动变形响应，
而强化学习策略决定了主动偏转输入，
二者的协同优化产生了远超纯主动或纯被动控制的气动性能提升，
为新一代智能变体飞行器设计提供了理论和方法基础
\cite{nairBioinspiredVariablestiffnessFlaps2023}。

上述研究表明，深度强化学习正在深刻重塑主动流动控制的研究范式，
从二维到三维、从单物理场到流固耦合、从层流到硬湍流，
DRL展现出了对复杂流动现象进行感知、建模和控制的系统级能力。
其核心优势在于能够绕过对Navier-Stokes方程的显式降阶建模，
直接从高维流场数据中提取与任务目标相关的控制律。
然而，当前AFC领域的研究仍多集中于中低雷诺数（$Re < 10^4$）
和几何外形相对简单的基准算例，
向真实飞行器尺度和复杂几何外形的工程化拓展仍面临
样本效率、策略泛化性和实时部署等技术挑战，
这些挑战同时也是本领域未来最具价值的研究方向。
\section{本文研究内容与结构安排}

基于上述研究背景与现状分析，本文以活性物质物理理论为基础，
将无人机滑翔问题建模为自推进体在非定常流场中的点对点迁移问题，
采用深度强化学习算法探索最小能耗迁移策略。
本文的核心研究目标是：（1）建立自推进体在非定常涡流场中的
运动学和能耗模型；（2）设计适用于高维连续动作空间的
深度强化学习框架；（3）揭示智能体在湍流环境中
实现低能耗导航的物理机制。本文的主要研究内容与结构安排如下：

第一章为绪论，阐述研究背景与意义，
从无人机滑翔、活性物质非平衡输运和强化学习流体控制
三个维度系统综述国内外研究进展，
明确本文的研究目标、技术路线与内容组织。

第二章建立自推进体在非定常流场中的运动学模型，
介绍双涡流流场的数学描述与拓扑结构特征，
推导自推进力的力学表达与推进功率、能量消耗的定量定义，
并给出点对点迁移任务的形式化描述与评价指标。

第三章详细阐述强化学习算法的设计与实现，
包括Soft Actor-Critic (SAC)算法的基本原理
与其在连续控制问题中的优势分析，
状态空间与动作空间的物理构造方法，
奖励函数的多目标设计策略与权重整定，
以及神经网络架构、训练流程与关键超参数的配置方案。

第四章展示仿真结果并进行深入的物理机制分析，
涵盖流场拓扑特征与智能体迁移轨迹的可视化对比、
能耗累积曲线的时间演化分析、不同自推进强度下
轨迹形态的系统性变化规律、训练过程的收敛性与稳定性评估、
随机起终点条件下的策略泛化能力测试，
以及基于速度-流场余弦相似度的“借流而行”物理机制的定量验证。

第五章对全文工作进行总结，提炼核心结论与理论创新点，
分析当前研究方法的适用边界与局限性，
并展望将该框架推广至三维流场、
多智能体协同导航和真实大气环境应用等未来研究方向。